%!TeX root=../tese.tex
\chapter{Calibrating an Analytical Cost Model for the Vision Pool}
\label{cap:aiconfigurator}

This chapter reports the construction and empirical validation of a
silicon-calibrated cost model for the vision stage of the disaggregated
serving architecture studied in this thesis. The model is built on
\textit{aiconfigurator} (AIC), NVIDIA's offline deployment planner, whose
performance database prices a model's operator graph against per-kernel
latencies collected on the target GPU. Three results are established.
First, the target consumer GPU (NVIDIA GeForce RTX~5090) was onboarded as a
new AIC system, and the two operator families that dominate vision-encoder
cost --- non-causal encoder attention and bf16 GEMM --- were collected on
its silicon under locked clocks. Second, the resulting analytical
prediction of the full vision-encoder forward pass was validated against
direct measurement of the same network on the same GPU at the same clocks:
the collected attention operator agrees with physical measurement within
3\%, and the end-to-end kernel-time ledger closes with no unexplained
residual. Third, the calibrated encoder model was composed with Dynamo's
discrete-event LLM-pool simulator in a first end-to-end two-pool
simulation, establishing the pipeline used for the autoscaling experiments
in the following chapters.

\section{Why silicon data is required}
\label{sec:aic-motivation}

AIC ships collected performance data for datacenter systems (H100, H200,
B200, B300, GB200, GB300), but its \texttt{encoder\_attention} operator ---
the non-causal, KV-cache-free attention used by ViT-style vision encoders
--- has no data for any workstation- or consumer-class GPU. This gap is
not softened by any fallback. Tracing the database's source-resolution
logic and exercising it empirically on a system without encoder data
(\texttt{rtx\_pro\_6000\_server}, the workstation sibling of the GPU used
here) shows that the SILICON, HYBRID, and EMPIRICAL query modes all raise a
hard \texttt{PerfDataNotAvailableError}; only the speed-of-light (SOL)
mode answers, and it underestimates the measured operator latency by a
factor of ${\approx}5.9$ at a representative shape (\SI{0.0061}{\milli\second}
predicted vs.\ \SI{0.0358}{\milli\second} measured at batch~1,
784~patches). The green entries that vision-language models show in AIC's
support matrix on such systems are explained by the harness supplying no
image dimensions: the encoder phase is skipped entirely, so the entry
certifies text-only costing. Image-aware costing of a multimodal model on
consumer hardware therefore requires collecting encoder data on that
hardware --- the work of this chapter.

\section{System onboarding and data collection}
\label{sec:aic-collection}

\subsection{Target system}

The vision pool's hardware class is represented by an NVIDIA GeForce
RTX~5090 (GB202 die, consumer Blackwell, compute capability~12.0,
\SI{32607}{\mebi\byte} GDDR7, observed board power limit \SI{600}{\watt},
driver 595.71.05, CUDA~13.2). The system was registered in AIC as
\texttt{rtx\_5090}: memory bandwidth and capacity follow the observed node
values; tensor-core throughput is scaled from the same-die
\texttt{rtx\_pro\_6000\_server} entry's per-SM, per-MHz throughput at the
official \SI{2407}{\mega\hertz} boost clock, yielding
\SI{419}{\tera\flop\per\second} dense bf16, consistent with published
specifications. All estimated fields are marked as such in the system
file for later refinement.

\subsection{Collection methodology}

Collection followed the upstream \texttt{aic-auto-collect} workflow, with
one audit artifact (script plus captured output) per phase. Three
methodological points matter for the validity of the data.

\paragraph{Kernel fidelity.} The vLLM encoder-attention collector does not
hard-code a kernel choice; it calls vLLM's own backend selector
(\texttt{get\_vit\_attn\_backend}) and benchmarks whatever kernel vLLM
would use in production on the target GPU. On compute capability 12.0
under vLLM~0.22.0 this resolves to FlashAttention
(\texttt{kernel\_source=vllm\_vit\_flash\_attn}), not the Triton fallback
that older documentation suggests for non-datacenter parts. The measured
kernel is therefore the deployed kernel by construction. Collection ran
inside the pinned production container (\texttt{vllm/vllm-openai:v0.22.0});
the collector module's compatibility gate (\texttt{vllm>=0.21.0}) rules
out the older framework versions that AIC's manifest still lists.

\paragraph{Clock control.} Consumer GPUs boost opportunistically, and
isolated microbenchmark kernels ride higher clocks than sustained serving
load. An initial collection at unlocked clocks produced latencies
${\approx}17\%$ faster than the final dataset. All data reported here was
recollected with the SM clock locked to \SI{2407}{\mega\hertz} and the
memory clock to \SI{14001}{\mega\hertz}
(\texttt{nvidia-smi -lgc/-lmc}, persistence mode enabled), matching the
clock at which the system file's throughput figures are stated. Notably,
the measured-to-predicted \emph{ratios} of Section~\ref{sec:aic-validation}
are identical between the unlocked and locked datasets --- the clock state
scales both sides equally and cancels --- but only the locked dataset
makes the absolute numbers interpretable.

\paragraph{Documented gaps.} On a \SI{32}{\gibi\byte} card the GEMM base
sweep's giant tail (e.g.\ $M{=}32768$, $N{=}65536$) cannot allocate; 1{,}157
such cases fail with out-of-memory errors and 276 were pre-skipped by a
\SI{4}{\gibi\byte} tensor-footprint cutoff, all recorded as gaps rather
than filled synthetically. Repeated allocation failures were additionally
found to corrupt CUDA-graph state and crash the collection process; the
production runs therefore pre-skip doomed shapes, clear the allocator
cache after failures, and restart the container between chunks. Seventeen
small shapes ($M{\le}2$) that spuriously reported out-of-memory in the
fragmentation wake of giant-shape failures were requeued in a fresh
container and all collected cleanly --- a failure mode worth noting for
any collection campaign on memory-constrained cards.

\subsection{Collected dataset}

The final dataset comprises \num{7679} encoder-attention cases (batch
$\in\{1,\dots,64\}$, sequence length $\in[1, 65536]$, 35 head
configurations including the $16\times 72$ geometry of the Qwen3-VL vision
tower; zero errors) and \num{34326} unique bf16 GEMM rows, both under
\texttt{systems/data/rtx\_5090/vllm/0.22.0/}. The dataset, system file,
and registration were contributed upstream as a draft pull request
(ai-dynamo/aiconfigurator \#1312).

\section{Validation against physical measurement}
\label{sec:aic-validation}

\subsection{Setup}

The validation target is the vision encoder of
\texttt{Qwen/Qwen3-VL-8B-Instruct}: a 27-block ViT (hidden size 1152, 16
heads of dimension 72, GELU MLP of width 4304) followed by a PatchMerger
projector into the language model's \num{4096}-dimensional embedding
space, plus three DeepStack mergers that project intermediate activations
from blocks 8, 16, and 24. The Qwen3-VL family shares this tower across
its 8B, 32B, and 235B variants, so results transfer to the larger models
unchanged. The tower (576M parameters, \SI{1.15}{\giga\byte} in bf16) was
extracted from the full checkpoint into a standalone
\texttt{Qwen3VLVisionModel} and verified to load and run independently;
its configuration is field-by-field identical to the vision configuration
AIC's model graph is built from, so prediction and measurement describe
the same network.

Measurement uses CUDA events around the full forward pass (patch input
resident on GPU), bf16, with 5 independent replications of 50 timed
iterations each (10 warm-up passes per replication); the SM clock is
sampled after every replication. The AIC prediction prices the model's
encoder operator graph --- per block: fused add-norm, QKV GEMM
($1152\to3456$), encoder attention (with rotary factor 0.5), output
projection ($1152\to1152$), MLP GEMMs ($1152\to4304\to1152$) and
activation; plus four projector instances (the main merger and the three
DeepStack mergers, $4608\to4608\to4096$ each) --- against the collected
silicon tables. Tensor-parallel all-reduce terms are zero at TP=1.

\subsection{Results}

Table~\ref{tab:aic-validation} compares measured wall-clock latency with
the AIC prediction at three image sizes. Confidence intervals are 95\%
$t$-intervals across the five replications.

\begin{table}[htb]
\centering
\caption{Vision-encoder forward latency: measurement vs.\ AIC prediction
(RTX~5090, bf16, batch 1, SM clock locked to \SI{2407}{\mega\hertz}).}
\label{tab:aic-validation}
\begin{tabular}{lrrrr}
\toprule
Image size & Patches & Measured (ms) & Predicted (ms) & Ratio \\
\midrule
$448^2$  & 784   & $11.360 \pm 0.006$  & 7.822  & 1.45 \\
$896^2$  & 3136  & $36.421 \pm 0.010$  & 31.973 & 1.14 \\
$1344^2$ & 7056  & $103.824 \pm 0.022$ & 89.347 & 1.16 \\
\bottomrule
\end{tabular}
\end{table}

The headline result is that \textbf{the analytical model matched the
empirical test}: a kernel-level profile of the measured forward pass
accounts for the entire gap, term by term, with no unexplained residual.
At $448^2$, the \SI{11.360}{\milli\second} wall time decomposes into
\SI{9.624}{\milli\second} of pure CUDA kernel time plus
\SI{1.736}{\milli\second} (15.3\%) of inter-kernel dispatch gaps --- a
property of the HuggingFace eager execution stack used for measurement,
deliberately outside AIC's scope. Within the kernel time, the attention
kernels sum to \SI{1.300}{\milli\second} against a predicted
\SI{1.342}{\milli\second} --- \textbf{agreement within 3\%} on the
operator that was actually collected. The GEMM total runs ${\approx}6\%$
above prediction (HuggingFace selects different cutlass variants than
vLLM's GEMM path), unfused elementwise operations and the projectors'
layer norms account for a further ${\approx}0.25$ ms, and two modules sit
outside AIC's operator graph altogether: the patch-embedding convolution
(\SI{0.459}{\milli\second}, 4.8\% of kernel time --- an inefficient
convolution shape, considerably more expensive than its FLOP count
suggests) and the DeepStack gather/concatenation
(\SI{0.248}{\milli\second}). Removing the two unmodeled modules leaves
\SI{8.92}{\milli\second} of AIC-comparable work against
\SI{7.82}{\milli\second} predicted, within ${\approx}12\%$, fully
attributed to the kernel-variant and fusion differences above. The
DeepStack merger MLPs themselves \emph{are} modeled: AIC instantiates the
projector four times, and forcing the instance count from four to one
scales the predicted projector cost by exactly $4.00\times$.

The gap's size ordering is consistent with this decomposition: it shrinks
from $1.45\times$ at $448^2$ to $1.14$--$1.16\times$ at larger images as
the fixed per-layer dispatch overhead amortizes over longer kernels.
The residual effects --- dispatch patterns, kernel fusion, the serving
framework's execution discipline --- are precisely the serving-stack
phenomena that the discrete-event simulation of the next phase models
explicitly; at this phase the analytical model's job is per-operator
kernel cost, and at that job it is validated.

\section{Two-pool composition}
\label{sec:aic-twopool}

The calibrated encoder model was composed with Dynamo's offline
LLM-pool simulator (\texttt{dynamo.replay}) in a minimal end-to-end test
of the thesis's two-pool architecture. Stage one simulates the vision
pool as a FIFO queue over $C$ encoder workers whose per-image service
time (\SI{7.82}{\milli\second} at $448^2$) comes from the AIC prediction
on the collected RTX~5090 silicon; each request's encoder delay (queue
wait plus service) is folded into its arrival timestamp at the language
pool. Stage two replays the composed trace --- input length equal to
text tokens plus the 196 post-merge vision tokens per image, with
disjoint prefix-hash sets so no spurious cross-image cache reuse occurs
--- through \texttt{dynamo.replay} in offline mode with a KV-aware router
over two simulated workers. A 50-request smoke workload (one $448^2$
image and 256 text tokens per request, \SI{20}{\milli\second} arrival
interval) produces a complete AIPerf-style report (mean time-to-first-token
\SI{32.4}{\milli\second}; 27.6 requests per second), confirming that
measured silicon, analytical encoder costing, queueing composition, and
the LLM-pool simulator interoperate end to end.

The first of the fidelity upgrades this pipeline invites --- AIC engine
timing for the language pool on the datacenter \texttt{h100\_sxm} system,
replacing the simulator's default polynomial timing --- has since been
carried out, driven by the real Azure multimodal trace rather than the
synthetic smoke workload. It required no engine rebuild: the shipped
inference runtime already exposes the engine-timing AIC surface, and
enabling it only meant co-locating the AIC performance model with the
replay runtime in one interpreter. On an identical trace window and
topology, the calibrated timing exposes markedly heavier tail latency than
the polynomial default (99th-percentile time-to-first-token
${\approx}1.2\times$, tail maxima $3$--$5\times$, per-user decode throughput
${\approx}12\%$ lower) --- precisely the burst-time behaviour a cost-blind
timing model flattens. Realistic image-size distributions, load sweeps to
each pool's saturation knee, and planner-in-the-loop simulation remain for
the phase that follows.

\section{Summary}

The vision-pool cost model now rests on measured silicon: \num{7679}
encoder-attention and \num{34326} GEMM data points collected on the
target consumer GPU under locked clocks, contributed upstream, and
validated against physical measurement of the real network on the same
hardware --- within 3\% on the collected attention operator, with every
remaining difference enumerated and bounded. The chapter also documents
two findings of independent value: AIC provides no usable fallback for
encoder timing on systems without collected data (its support matrix can
mask this by skipping the encoder phase), and unlocked consumer-GPU
clocks silently inflate collected performance by ${\approx}17\%$ while
leaving validation ratios unchanged. The composed two-pool simulation
closes the phase: the pipeline from silicon to serving-level metrics is
in place for the autoscaling experiments that follow.
